\documentclass[]{informeutn}

\usepackage{circuitikz}
\usepackage{colortbl}
\usepackage{multirow}

% Datos del informe
\materia{Dispositivos Electrónicos}
\titulo{Trabajo Práctico 1}
\autores{Cabaleiro Martin\\Cortesini Luciano\\Ernst Pedro}
\fecha{8 / 4 / 2025}

\begin{document}

\maketitle

\tableofcontents
\setcounter{page}{1}
\thispagestyle{plain}

\chapter{Introducción}

En este trabajo práctico de laboratorio se estudiará el comportamiento de un circuito resistivo frente a señales continuas y alternas.

Para ello, en primer lugar se plantea el circuito y se analiza de forma analítica utilizando las leyes de Kirchhoff (ley de nodos y
ley de mallas), la ley de Ohm y el cálculo de resistencias equivalentes. Luego, se realizará una simulación del circuito mediante el
software LTSpice con el objetivo de verificar los resultados obtenidos teóricamente. Finalmente, el circuito se implementará
físicamente sobre una placa protoboard, y se llevarán a cabo mediciones de tensiones y corrientes. Para las señales continuas se
utilizará un multímetro, mientras que para las señales alternas se utilizará un osciloscopio.

El objetivo principal de este trabajo es comparar los resultados teóricos, simulados y experimentales, analizar el comportamiento 
del circuito resistivo ante diferentes tipos de señal, y practicar el uso de los instrumentos y elementos disponibles en el laboratorio.

\chapter{Análisis del circuito}

\begin{center}
  \begin{circuitikz} [scale=1, transform shape]
    \draw  (0,0) to[battery, invert, l=$Vs$] (0,3)
        to[R, l=$R_1$,] (3,3) node[circ]{}
        to[R, l=$R_2$] (3,0);
    \draw (3,3) to[short] (5,3)
        to[R, l=$R_3$] (5,0)
        to[short] (3,0) node[circ]{}
        to[short] (0,0);
    \draw[->, thin] (1,2.5) -- (2,2.5) node[pos=0.5, below] {$I_1$};
    \draw[->, thin] (2.5,2) -- (2.5,1) node[pos=0.5, left] {$I_2$};
    \draw[->, thin] (4.5,2) -- (4.5,1) node[pos=0.5, left] {$I_3$};
  \end{circuitikz}
\end{center}

\chapter{Implementación y mediciones}

\begin{table}[H]
  \centering
  \renewcommand{\arraystretch}{1.7}
  \begin{tabular}{|c|c|c|c|c|c|}
    \hline
    \multicolumn{2}{|c|}{} & \textbf{$V_s$} & \textbf{$R_1$} & \textbf{$R_2$} & \textbf{$R_3$} \\
    \hline
    \multirow{3}{*}{Tension} & Análisis   &                 &                &                &                \\
    \cline{2-6}
                             & Simulación &                 &                &                &                \\
    \cline{2-6}
                             & Medición   &                 &                &                &                \\
    \hline
    \multirow{3}{*}{Corriente} & Análisis   &               &                &                &                \\
    \cline{2-6}
                               & Simulación &               &                &                &                \\
    \cline{2-6}
                               & Medición   &               &                &                &                \\
    \hline
  \end{tabular}
\end{table}

\chapter{Conclusión}

\section{Rúbrica}

\begin{table}[H]
  \centering
  \renewcommand{\arraystretch}{1.7}
  \begin{tabular}{|c|c|c|}
    \hline
   \rowcolor[gray]{0.8}
   \textbf{Tarea} & \textbf{Puntuación} & \textbf{Máximo} \\
   \hline
   Presentación del informe & & 30 \% \\
   \hline
   Explicación del proceso de medición & & 40 \% \\
   \hline
   Defensa de las conclusiones & & 30 \% \\
   \hline
   \textbf{Total} & & \textbf{100 \%} \\
   \hline
  \end{tabular}
\end{table}

\end{document}

